\chapter*{Εκτεταμένη Ελληνική Περίληψη}

\begin{center}
    {\Large \textbf{Σχεδίαση και ανάπτυξη φορητής εφαρμογής παρακολούθησης διάθεσης}}
\end{center}

\vspace{5mm}

\begin{minipage}[t]{0.45\textwidth}
    \begin{center}
        ΟΝΟΜΑΤΕΠΩΝΥΜΟ ΦΟΙΤΗΤΗ \\[2em]
        Αλέξανδρος Τσαπάρας \\[0.5em]
    \end{center}
\end{minipage}%
\hfill
\begin{minipage}[t]{0.45\textwidth}
    \begin{center}
        ΟΝΟΜΑΤΕΠΩΝΥΜΟ ΕΠΙΒΛΕΠΟΝΤΟΣ \\[0.2em]
        Χρήστος Σιντόρης \\[0.5em]
    \end{center}
\end{minipage}

\vspace{3em}

\noindent Η παρούσα διπλωματική εργασία εστιάζει στη σχεδίαση και ανάπτυξη μιας φορητής εφαρμογής παρακολούθησης διάθεσης με στόχο την ενίσχυση της ψυχικής υγείας των φοιτητών. Οι φοιτητές αντιμετωπίζουν συχνά έντονη ψυχολογική πίεση, κυρίως λόγω των ακαδημαϊκών και κοινωνικών προσδοκιών. Η εφαρμογή που αναπτύχθηκε έχει ως στόχο να βοηθήσει τους φοιτητές να καταγράφουν τη διάθεσή τους καθημερινά, να αναλύουν συναισθηματικά μοτίβα και να λαμβάνουν συμβουλές για την αντιμετώπιση των προβλημάτων τους, με απλό και εύχρηστο τρόπο.\vspace{5mm} \\
Η ανάπτυξη της εφαρμογής βασίστηκε σε τρεις κύριες φάσεις: (1) η ανασκόπηση της βιβλιογραφίας και ανάλυση των υφιστάμενων λύσεων παρακολούθησης διάθεσης, (2) η ανάπτυξη και σχεδίαση της εφαρμογής, και (3) η αξιολόγηση της εφαρμογής μέσω δοκιμών χρήσης από φοιτητές.

\section*{Κίνητρο και Στόχος}
Το κίνητρο για την ανάπτυξη της εφαρμογής προέκυψε από την προσωπική μου εμπειρία ως φοιτητής, καθώς και από την ανάγκη αντιμετώπισης της ψυχολογικής πίεσης που βιώνουν πολλοί φοιτητές. Πολλοί φοιτητές αποφεύγουν να ζητήσουν βοήθεια, κυρίως λόγω του κοινωνικού στίγματος ή της έλλειψης ενημέρωσης για τους διαθέσιμους πόρους. Η εφαρμογή επιτρέπει την παρακολούθηση της διάθεσης με τρόπο που ενισχύει την αυτογνωσία και παρέχει πρακτικά εργαλεία για τη βελτίωση της ψυχικής ανθεκτικότητας.\vspace{5mm} \\
Ο κύριος στόχος της διπλωματικής ήταν να δημιουργηθεί μια εφαρμογή που:
\begin{itemize}
    \item Να επιτρέπει τη συστηματική καταγραφή της διάθεσης των φοιτητών.
    \item Να παρέχει ανάλυση συναισθηματικών μοτίβων μέσω γραφημάτων.
    \item Να προσφέρει εξατομικευμένες συμβουλές και υποστήριξη.
    \item Να ενισχύσει τη συνολική ψυχική υγεία και ευεξία των φοιτητών.
\end{itemize}

\section*{Μεθοδολογία και Σχεδίαση}
Η ανάπτυξη της εφαρμογής ακολούθησε μια δομημένη μεθοδολογία που περιλάμβανε:
\subsection*{Ανασκόπηση Βιβλιογραφίας}
Μέσω της ανασκόπησης της βιβλιογραφίας, αναλύθηκαν έρευνες σχετικά με την ψυχική υγεία των φοιτητών και τις επιπτώσεις της διάθεσης στην ακαδημαϊκή τους απόδοση. Συγκεκριμένα:
\begin{itemize}
    \item Έρευνες από το University of Applied Sciences, Northern Netherlands, και το Renmin University of China ανέδειξαν τις επιπτώσεις του άγχους και της ψυχολογικής πίεσης.
    \item Μελέτες σε ιδιωτικά πανεπιστήμια της Ινδονησίας έδειξαν πως η αρνητική διάθεση επηρεάζει σημαντικά τη μαθησιακή διαδικασία.
\end{itemize}

\subsection*{Ανάλυση Απαιτήσεων (PACT Model)}
Η ανάλυση των αναγκών των χρηστών έγινε μέσω του μοντέλου PACT, το οποίο αναλύει:
\begin{itemize}
    \item \textbf{People (Χρήστες)}: Η εφαρμογή απευθύνεται σε φοιτητές με διαφορετικές ανάγκες και επίπεδα εμπειρίας.
    \item \textbf{Activities (Δραστηριότητες)}: Η καθημερινή καταγραφή διάθεσης και η ανάλυση δεδομένων είναι οι κύριες δραστηριότητες.
    \item \textbf{Context (Πλαίσιο)}: Χρήση σε διάφορα περιβάλλοντα, όπως πανεπιστήμια ή σπίτι.
    \item \textbf{Technology (Τεχνολογία)}: Χρήση φορητών συσκευών με ασφαλή διαχείριση δεδομένων μέσω cloud.
\end{itemize}

\subsection*{Προσωπικότητες (Personas) και Σενάρια Χρήσης}
Για την καλύτερη κατανόηση των χρηστών, δημιουργήθηκαν τρεις διαφορετικές προσωπικότητες:
\begin{itemize}
    \item \textbf{Δήμητρα}: Φοιτήτρια με υψηλές ακαδημαϊκές φιλοδοξίες που αντιμετωπίζει άγχος.
    \item \textbf{Ιωάννης}: Φοιτητής που προσπαθεί να διατηρήσει την ισορροπία μεταξύ ακαδημαϊκής και προσωπικής ζωής.
    \item \textbf{Κώστας}: Φοιτητής που βιώνει προσωπικές δυσκολίες και αισθάνεται απομονωμένος.
\end{itemize}

\subsection*{Σχεδίαση Διεπαφής Χρήστη}
Η διεπαφή της εφαρμογής σχεδιάστηκε για να είναι απλή και εύχρηστη. Κάποιες από τις βασικές λειτουργίες περιλαμβάνουν:
\begin{itemize}
    \item Καταγραφή διάθεσης με τη χρήση emojis.
    \item Απεικόνιση συναισθηματικών μοτίβων μέσω γραφημάτων.
    \item Εξατομικευμένα μηνύματα υποστήριξης μετά την καταγραφή της διάθεσης.
\end{itemize}

\section*{Ανάπτυξη της Εφαρμογής}
\subsection*{Front-end Ανάπτυξη}
Η ανάπτυξη του front-end της εφαρμογής έγινε με τη χρήση της πλατφόρμας React Native και Expo, επιτρέποντας την πολυπλατφορμική ανάπτυξη για iOS και Android.

\subsection*{Back-end Ανάπτυξη}
Το back-end αναπτύχθηκε με χρήση Node.js και Express, διασφαλίζοντας την ταχύτητα και αποδοτικότητα. Η βάση δεδομένων σχεδιάστηκε με PostgreSQL και φιλοξενείται στην πλατφόρμα Xata, η οποία προσφέρει ασφάλεια και κλιμακούμενη διαχείριση δεδομένων.

\section*{Αξιολόγηση και Αποτελέσματα}
Η αξιολόγηση της εφαρμογής πραγματοποιήθηκε μέσω δοκιμών χρήσης με φοιτητές. Οι χρήστες βρήκαν την εφαρμογή εύχρηστη και αναγνώρισαν τη χρησιμότητά της στην παρακολούθηση της διάθεσής τους. Επιπλέον, αναφέρθηκαν προτάσεις για βελτίωση των γραφημάτων και της εξατομίκευσης της εμπειρίας.

\section*{Συμπεράσματα}
Η εφαρμογή παρακολούθησης διάθεσης επιτυγχάνει τον στόχο της, προσφέροντας στους φοιτητές ένα χρήσιμο εργαλείο για τη βελτίωση της ψυχικής τους υγείας. Παρά τις μικρές βελτιώσεις που εντοπίστηκαν, η εφαρμογή παρέχει μια αξιόπιστη λύση για την καταγραφή και ανάλυση διάθεσης.

\section*{Μελλοντικά Βήματα}
Τα επόμενα βήματα περιλαμβάνουν τη βελτίωση της διεπαφής χρήστη, την προσθήκη νέων χαρακτηριστικών όπως προτάθηκαν από τους χρήστες, και την ανάπτυξη της εφαρμογής για διάθεση στο App Store και Google Play.
